
%%%%%%%%%%%%%%%%%%%%%%%%%%%%%%%%%%%%%%%%%%%%%%%%%%%%%%%%%%%%%
\section{模型检验与分析}

为验证本文所建模型的有效性与稳健性，本章从交叉验证、残差分析与灵敏度分析三个角度进行检验。

\subsection{交叉验证检验}

本文模型的核心结论均建立在解析推导之上，因此采用"解析结论与独立数值求解逐点比对"的方式做交叉验证，比对结果汇总于表 \ref{tab:cv}。

\begin{table}[H]
	\centering
	\caption{解析结论与独立数值求解的交叉验证}
	\label{tab:cv}
	\begin{tabularx}{\textwidth}{p{2.4cm}p{3.6cm}CC}
		\toprule
		所属问题 & 检验对象 & 检验方式 & 最大相对偏差 \\
		\midrule
		问题一 & Thales 覆盖判据 & 与逐点距离判据比对（4860 组） & $0$（结论完全一致） \\
		问题一 & 旋转卡壳直径 & 与顶点对暴力枚举比对 & $<10^{-9}$ m \\
		问题一 & 覆盖率分组统计 & 按交会角九区间交叉核对 & 计数完全一致 \\
		问题二 & 最优基线 $b^{*}=\sqrt2r$ & 与一维精细搜索比对（8 个检验点） & $<10^{-5}$ \\
		问题二 & 最优交会角 $54.74^\circ$ & 与数值最优角比对 & $<10^{-2}$ 度 \\
		问题二 & 候选区域边界 & 与网格扫描次水平集比对 & 一致（步长内） \\
		问题三 & 观测点数下界 $N_q\ge7$ & Kershner 与数值优化双向验证 & 两法一致 \\
		问题三 & 覆盖网覆盖率 & $1.5\times10^6$ 点稠密采样 & 覆盖率 $99.999867\%$ \\
		问题三 & 时间模型分项 & 与逐次动作累加核对 & $0$ \\
		问题四 & 背向不可测定理 & 4000 组随机构造验证 & 反例 $0$ 组 \\
		问题四 & 零漏测门槛 $180^\circ$ & 与网络规模权衡曲线核对 & 一致 \\
		\bottomrule
	\end{tabularx}
\end{table}

如表 \ref{tab:cv} 所示，全部检验项的偏差均处于数值精度范围内，未发现解析结论与数值结果之间的实质性分歧。特别值得注意的是问题三中观测点数下界的双向验证：一方面 Kershner 定理从理论上说明 $6$ 个半径 $1000$ 米的圆盘最多覆盖半径 $1732.05$ 米的圆域，另一方面对纯 $6$ 环点构型做覆盖半径优化得到的数值结果为 $1039.23$ 米，二者从不同角度共同确认了 $7$ 是必需的观测点数。

\subsection{残差分析检验}

本文模型的主要误差来源是示向度的角度误差，因此以"定位误差随源距离与基线长度的分布"作为残差分析的载体。在问题二中，把式 $E(b;r)=\frac{\delta}{b}\sqrt{(r^2+b^2)(4r^2+b^2)}$ 在固定 $r$ 下关于 $b$ 求得的解析曲线与在可行域上直接采样的数值期望逐点比较，二者在 $b\in[20,3200]$ 米范围内完全重合，残差随 $b$ 增大而单调衰减，呈现典型的模型误差形态而非结构性偏差。

在问题一中，对 $4860$ 个有效构型统计"最小包围圆半径与直径之比"这一残差型指标，其均值为 $0.500001$、最大值仅为 $0.500280$，紧贴理论下界 $0.5$ 而远低于 Jung 上界 $0.577350$。该结果表明定位区域的形状高度集中于"瘦长"形态，模型对临界情形的刻画是准确的；若该比值在统计上显著偏离 $0.5$，则说明交会几何存在系统性的形状偏好，此时 Thales 判据的适用性就需要重新审视。

在问题三中，时间模型的各分项之和与逐次动作累加结果完全一致，残差为零，说明虚拟时间口径的换算没有遗漏任何计费项目。在问题四中，对"七点网在半径 $400\sim800$ 米环带内保证不漏测"这一结论做分层残差检查，该层最大方位间隔为 $176.48^\circ$，与门槛 $180^\circ$ 之间的余量为 $3.52^\circ$，余量偏小，提示该环带处于临界状态，实际部署时应留出额外的观测密度裕量。

\subsection{灵敏度分析}

灵敏度分析围绕四个关键参数展开，结果汇总于表 \ref{tab:sens}。

\begin{table}[H]
	\centering
	\caption{关键参数的灵敏度分析汇总}
	\label{tab:sens}
	\begin{tabularx}{\textwidth}{LLL}
		\toprule
		参数 & 变化范围 & 主要影响 \\
		\midrule
		示向度误差上界 $\delta$ & $0.25^\circ\sim4^\circ$ & 问题一直径圆覆盖率由 $99.93\%$ 降至 $96.39\%$；\newline 问题三逼近阈值 $r_{\mathrm{thr}}$ 由 $1528$ m 降至 $95$ m \\
		侧移系数 $\beta$ & $0\sim0.5$ & 问题三平均定位清除时间由 $297$ s 升至 $366$ s \\
		机动成本权重 $\lambda$ & $0\sim1.0$ & 问题二最优基线由 $1401$ m 收缩至 $461$ m，\newline 方位角向 $54.74^\circ$ 靠拢 \\
		外环半径 $\rho_o$ & $1500\sim2200$ m & 边界源可见窗口半角由 $33.75^\circ$ 收窄至 $26.63^\circ$ \\
		\bottomrule
	\end{tabularx}
\end{table}

如表 \ref{tab:sens} 所示，第一，示向度精度是最敏感的杠杆。问题一中当 $\delta$ 由 $1^\circ$ 收紧到 $0.25^\circ$ 时，直径圆覆盖率由 $98.99\%$ 升至 $99.93\%$，失败构型的最大超出量由 $1.17$ 米降至 $0.05$ 米；问题三中逼近阈值与 $\delta$ 成反比，$\delta$ 每减半所需的逼近距离翻倍，意味着定位阶段可以更早停止逼近，时间代价显著下降。第二，侧移系数 $\beta$ 对总时间近似线性，是其直接控制变量，建议通过现场试验标定。第三，机动成本权重 $\lambda$ 揭示了问题二与问题三之间的耦合：当移动代价不可忽略时，最优基线向"短基线加最优交会角"的方向漂移。第四，外环半径的可行窗口较窄，$\rho_o$ 过大将导致边界源的可见窗口收窄，反而需要更密的环点，故外环半径应取在 1900 米附近。
